\begin{table}[htbp]
\centering
\footnotesize
\caption{\textbf{Mean peaker operating cost, by scenario and year.} Across the EU-27, the United Kingdom and Switzerland; short-run marginal cost in EUR/MWh-e, fuel and carbon only, no capital. Hydrogen's cost falls as delivered hydrogen cheapens and the turbine improves; the carbon-priced gas peaker rises with the carbon path (about 70 to 350 EUR/tCO$_2$ in the two high-carbon scenarios). The two cross only late and only under H2 Push (221 versus 239 in 2050 on this MEAN basis; on the median the same year is a tie at 238 against 239, and the AE submission reports the median because this distribution is bimodal and the mean is dragged by the salt-cavern tail, so the sign of the crossing depends on which statistic is quoted); in the lower-carbon scenarios gas stays cheaper throughout. The EU-27+UK+CH mean understates the salt-cavern markets, where the hydrogen cost is far lower (Figure~\ref{fig:crossover}).}
\label{tab:opcost}
\begin{tabular}{llrrrr}
\toprule
Scenario & Peaker & 2025 & 2030 & 2040 & 2050 \\
\midrule
Current Policies & Gas & 114 & 121 & 141 & 143 \\
& H$_2$ & 667 & 605 & 487 & 416 \\
\addlinespace
Stated Policies & Gas & 127 & 147 & 165 & 167 \\
& H$_2$ & 667 & 417 & 316 & 281 \\
\addlinespace
Net Zero & Gas & 135 & 186 & 224 & 239 \\
& H$_2$ & 667 & 417 & 316 & 281 \\
\addlinespace
H2 Push & Gas & 135 & 186 & 224 & 239 \\
& H$_2$ & 667 & 393 & 267 & 221 \\
\addlinespace
\bottomrule
\end{tabular}
\\[8pt]{\footnotesize\textit{Source: Authors' contribution.}}
\end{table}
